El objetivo de esta tesina era desarrollar un simulador de dinámica vehicular capaz de ejecutarse en tiempo real con el conductor dentro del lazo. Las secciones que siguen resumen lo que el trabajo permite afirmar sobre ese objetivo, sobre las herramientas empleadas, sobre los límites que se encontraron y sobre mejoras posibles al proyecto.

\section{Conclusiones}

\subsection*{Es viable construir un simulador de conducción en Modelica}

La conclusión principal es positiva. Se pudo construir una biblioteca de dinámica vehicular multicuerpo en Modelica que acepta comandos de un conductor humano, o una herramienta externa, avanza la simulación al ritmo del tiempo real y devuelve la pose del vehículo a un entorno gráfico, todo con una latencia aceptable para la interacción con el ser humano. La mediana entre la recepción del comando y el envío de la pose es de entre $13{,}1$ y $15{,}6$ ms  y el percentil $99$ no supera los $31$ ms.



% El resultado no era evidente de antemano. Modelica es un lenguaje pensado para el modelado declarativo y la simulación por lotes, y su ecosistema no contempla la ejecución interactiva: no existe en él una noción de plazo ni de presupuesto de cómputo por paso, y el propio compilador asume que el integrador puede tomarse el tiempo que necesite. Llevar un modelo de esta clase a un lazo cerrado con un conductor exigió, por lo tanto, sacar el muestreo y el acompasamiento fuera del sistema de ecuaciones y dejar dentro únicamente la física continua. Esa separación, adoptada inicialmente para sortear una limitación de la herramienta, resultó ser la decisión de arquitectura más acertada del trabajo.

\subsection*{La composición orientada a objetos como herramienta de diseño}

Modelica organiza los modelos como componentes conectados por interfaces, esa estructura resultó determinante a la hora de implementar la biblioteca. Por ejemplo, las tres topologías de suspensión implementadas comparten una misma interfaz, por lo que cambiar la geometría de suspensión de un vehículo es tan simple como cambiar un bloque y sus parámetros con otro basado en esa misma clase. Esto se aplica a cualquier componente del vehículo, lo que permite la creación de diferentes variantes de manera muy simple.

Esto además permite la fácil expansión de la librería, por ejemplo, las barras estabilizadoras fue un componente que fue agregado a último momento, y lo único que se tuvo que hacer para implementarlo fue crear el modelo que lo represente y conectarlo a los marcos ya expuestos.
% Dos observaciones de este trabajo muestran que la propiedad no es solo declarativa. La primera es que los tres vehículos presentan exactamente el mismo número de estados continuos, treinta, pese a montar suspensiones de doble triángulo, de puntal telescópico y multibrazo: la interfaz común no impone un costo de representación, porque el cierre analítico del bucle cinemático evita introducir restricciones adicionales. La segunda es que incorporar la barra estabilizadora, un componente que no existía cuando se diseñaron las suspensiones, no requirió modificar ninguno de los modelos previos: bastó con escribir un componente nuevo que se conecta a los marcos ya expuestos.

% La parametrización también cumple un papel similar. Dos de tres vehículos representados tienen la misma geometría de suspensión, sin embargo son dos tipos de autos muy diferentes ( 

% La parametrización cumple el mismo papel a otra escala. Los tres vehículos representados —un monoplaza de escuela, una berlina de tracción trasera y un compacto de tracción delantera— se obtienen del mismo código cambiando masas, rigideces, radios y puntos de anclaje. Construir un cuarto vehículo no exige programar, sino medirlo.

\subsection*{El desacoplamiento por protocolo estándar}

La comunicación entre el simulador, el programa que controla el dispositivo de entrada y el entorno de visualización se resolvió utilizando UDP con un formato binario estandarizado. Esta elección de usar un protocolo de red en vez de, por ejemplo, una comunicación directa entre procesos, permite sustituir cada una de las piezas sin tocar las demás.

Los tres programas mencionados anteriormente son tres procesos independientes que solo comparten el mismo formato de mensajes. Durante el desarrollo se cambió la forma de ejecutar el modelo (de un ejecutable de OpenModelica a un FMU controlado desde Python) sin modificar ninguna línea del visualizador ni del programa de entrada. Esto permite en un futuro cosas como cambiar el visualizador por uno más avanzado, o distribuir los procesos en máquinas separadas.


\subsection*{El desempeño es holgado con CVODE}

A la hora de elegir integradores hay una clara elección, CVODE es a la vez el más rápido y el más robusto de todos los ensayados. Es entre $1{,}1$ y $1{,}3$ veces más rápido que CVODE, a la vez que más robusto en los casos de prueba y $3$ a $22$ veces más rápido que \verb|gbode|.

Con esa elección el márgen con respecto al tiempo real es más que suficiente. Ejecutando el modelo en lazo abierto está entre $2{,}4$ y $16{,}1$ veces con respecto al tiempo real. Al cerrar el lazo el márgen en menor, lo que es esperado debido a que la entrada deja de ser constante y el sistema debe coordinar la comunicación, pero sigue siendo suficiente, la latencia total se mantiene en el orden de la decena de milisegundos, lo que deja espacio para agregar modelos de suspensión o neumáticos más complejos por ejemplo.



\subsection*{Lo que limita la fidelidad no es el cómputo}

Los ensayos de fidelidad muestran que el techo del simulador no está en la parte numérica. Con la tolerancia elegida el apartamiento de la trayectoria respecto a una solución de referencia es de entre $47$ y $275$ micrómetros al cabo de diez segundos de conducción, incluso con la tolerancia más permisiva ensayada el error máximo es de $2{,}3$ cm. Son cifras con órdenes de magnitud menores que la incertidumbre con la que se conocen las masas, rigideces y geometría de las suspensiones.

Las limitaciones encontradas fueron de el último tipo mencionado. El reparto de cargas resulta uniforme porque el chasis se modela como un cuerpo de masa homogénea, lo que afecta a algunos modelos y la relación del conjunto resorte-amortiguador del Civic queda fuera del rango real porque no se tienen sus puntos de anclaje reales. Estos problemas se pueden solucionar agregando otra forma de modelar el chasis en el primer caso, y con datos en el segundo, en ambos casos es sencillo adaptar la solución al sistema ya existente.


% La conclusión práctica es directa. Disponer de medidas precisas de las dimensiones y de los puntos de anclaje de un vehículo real rendiría, en fidelidad, mucho más que cualquier refinamiento adicional del integrador o de la ley de contacto. El trabajo lo comprobó en sentido inverso más de una vez: corregir la ubicación de un anclaje cambió la frecuencia propia de una suspensión, su curva de caída y su vía, sin tocar una sola ecuación.

\subsection*{Las conmutaciones son la dificultad numérica dominante}

Un patrón que se repitió en todos los experimentos, y en particular en todos los obstáculos encontrados es que la principal dificultad numérica observada se debe a una no linealidad conmutativa y no al tamaño del sistema ni a una complejidad cinemática. Por ejemplo, los topes de suspensión son el origen de las fallas de integración más persistentes y el escalón que introduce cada actualización del comando del conductor multiplica el costo con respecto a una entrada constante.

Por lo tanto, es más rentable invertir esfuerzo en suavizar las conmutaciones que en intentar ajustar parámetros del integrador. Vale la pena prestar atención a las discontinuidades que se presentan cerca de la condición de operación habitual.

% La lección, transferible a modelos de esta clase, es que el esfuerzo invertido en suavizar de forma controlada las conmutaciones rinde más que el invertido en ajustar los parámetros del integrador. Conviene además prestar atención especial a las discontinuidades situadas en la condición de operación habitual: una regularización mal ubicada, justo donde el sistema pasa la mayor parte del tiempo, puede costar más que todo el resto del modelo.

\subsection*{OpenModelica como herramienta}

El balance sobre la herramienta tiene altos y bajos. Como entorno de modelado y de ejecución en lazo abierto, OpenModelica funcionó sin problemas, permite compilar los modelos multicuerpo con múltiples integradores, modificar sus parámetros de forma sencilla y expone herramientas de diagnóstico y visualización que permiten comprender como funciona el sistema.

% El balance sobre la herramienta es de dos caras. Como entorno de modelado y de ensayo en lazo abierto, OpenModelica cumplió sin reservas: compila modelos multicuerpo de varios miles de ecuaciones, ofrece tres integradores con los que contrastar resultados, expone las herramientas de diagnóstico necesarias para entender por qué un sistema es difícil de integrar, y permitió construir y validar la totalidad de la biblioteca.

Como entorno de ejecución en tiempo real mostró un problema que llevó a tener que buscar alternativas para este tipo de entorno: un reloj periódico de muestreo impide integrar con CVODE.

% Como entorno de ejecución en tiempo real, en cambio, mostró límites concretos. Se identificaron y documentaron cuatro defectos, cada uno con su caso mínimo de reproducción, y dos de ellos afectan de lleno al lazo cerrado: un reloj periódico de muestreo impide integrar con CVODE, y los FMU de co-simulación terminan durante el avance. El primero es el que obligó a trasladar el muestreo fuera del modelo; el segundo, a exportar como FMU de intercambio de modelos y proveer el integrador desde la aplicación externa.

Esta limitación no tiene que ver con el lenguaje Modelica, si no, viene de su implementación concreta del integrador CVODE específicamente. Aquí el carácter de software libre de OpenModelica permite la modificación de esta implementación, lo que hace que si se quisiese solucionar esta dificultad, se puede hacer sin problemas, en cambio una herramienta propietaria haría que se tenga que buscar una alternativa de forma obligatoria.

% Corresponde señalar que ninguna de esas limitaciones es del lenguaje Modelica, sino de una implementación concreta y de una versión concreta. Y aquí el carácter de software libre de OpenModelica deja de ser un detalle de licencia para convertirse en una propiedad técnica del proyecto: los cuatro defectos están documentados de forma reproducible, el código que los origina es público, y nada impide —a este trabajo o a otro— corregirlos en el propio compilador en lugar de rodearlos. Una herramienta cerrada con los mismos defectos habría dejado como único camino el que aquí se siguió, que es evitarlos.
\section{Trabajo futuro}

El simulador desarrollado es funcional y posee una arquitectura que permite su crecimiento sin rehacer lo construido. A continuación hay algunos puntos en los que el simulador se podría mejorar, las cuales podrían ser desarrolladas a partir de la base establecida en esta tesina.

% El simulador desarrollado es funcional y su arquitectura admite crecer en varias direcciones sin rehacer lo construido. Las líneas que siguen se ordenan de menor a mayor alcance: primero las que consisten en agregar componentes sobre las interfaces ya definidas, luego las que exigen revisar un modelo físico, y por último las que involucran al entorno de ejecución.

\subsection{Ampliación del catálogo de componentes}

La extensión más directa es agregar modelos nuevos sobre las interfaces definidas. Cada paquete de componentes de la biblioteca se define a partir de una clase base que fija su interfaz, por lo que un modelo nuevo se puede crear a partir de escribir la física correspondiente y adhiriéndose a la interfaz ya escrita, lo que permite que sea agregado a un vehículo sin modificar el resto de elementos.

Por ejemplo, un diferencial autoblocante sería una buena adición al sistema. El diferencial actual reparte el par por partes iguales en las dos ruedas, en cambio, un autoblocante transfiere par hacia la rueda con más agarre entre ambas salidas. Implementarlo simplemente consiste en crear un modelo a partir de la interfaz de diferencial y escribir la ley sobre los elementos de esta. De esta misma forma se pueden agregar geometrías de suspensión adicionales, frenos con distribución regulable entre ejes, entre otros.


\subsection{Modelo de transmisión y motor}

El tren motriz es actualmente la parte más simplificada del vehículo: el par motor se representa mediante un bloque que entrega un valor proporcional al input del acelerador, sin curva de par, relaciones de caja o embrague. Es suficiente para acelerar el vehículo y estudiar su comportamiento dinámico, pero se requeriría un modelo más complejo para simular correctamente un vehículo real.

El punto de conexión ya existe hoy en día, es la brida de par que alimenta al diferencial, por lo que agregar esto al sistema simplemente consiste en agregar el modelo y un input adicional de selección de marchas.

% La cadena de tracción es, hoy, la parte más simplificada del vehículo: el par motor se representa mediante un bloque que entrega un valor proporcional al comando del acelerador, sin curva de par en función del régimen, sin relaciones de caja y sin embrague. Es suficiente para acelerar el vehículo y estudiar su comportamiento dinámico, pero no para reproducir la conducción de un automóvil real, en la que el punto de cambio, la caída de par entre marchas y la respuesta del motor forman parte central de lo que el conductor percibe.

% Un modelo apropiado incluiría la curva de par en función del régimen, una caja con relaciones seleccionables, el embrague y la inercia rotacional de la cadena. La biblioteca \textit{Rotational} de Modelica provee los elementos necesarios, y el punto de conexión ya existe: es la brida de par que hoy alimenta el diferencial. Esta línea tiene además una consecuencia para el conductor humano, porque incorpora una entrada más —el cambio de marcha— al conjunto de comandos que el lazo ya transporta.

\subsection{Terreno con propiedades de superficie}

El modelo de contacto actual toma el coeficiente de fricción como un parámetro del vehículo, por lo que toda la pista tiene el mismo agarre. Darle propiedades locales al terreno leyendo el coeficiente de fricción en el punto de cada rueda, en vez de leerlo de un parámetro permitiría representar zonas de distinto agarre, superficies mojadas o pasto fuera de la pista, y lo haría sin alterar la estructura del modelo de contacto.

A esta mejora se lo podría sumar otra de consistencia. El terreno existe dos veces, una en el modelo físico y otra en el entorno gráfico, si se actualiza una, la otra también tiene que modificarse para que coincida. Unificarlas evitaría inconsistencias entre lo que ve el conductor y lo que se está calculando en la simulación, es particularmente útil en circuitos donde la geometría no es trivial.


\subsection{Modelo de neumático}

El neumático es un componente que mejorarlo aumentaría la fidelidad de forma significativa. La implementación actual es una simplificación de lo que se puede encontrar en una herramienta comercial, el desarrollo de un modelo de neumático realista podría ser una tesina por sí sola. Una mejora posible, por ejemplo, sería tener un modelo que habilite la elección de compuestos diferentes de neumáticos, esto permitiría poder diferenciar entre ruedas de calle o de competición, agregando fidelidad dependiendo del tipo de auto modelado.


% El neumático es el componente cuya mejora más incidiría en la fidelidad, por dos motivos distintos. El primero es de parametrización: un modelo que distinga compuestos permitiría diferenciar un neumático de calle de uno de competición, que hoy se representan variando un único coeficiente de fricción. El segundo es estructural y de mayor alcance: la fuerza lateral se modela como viscosa, proporcional a la velocidad lateral de deslizamiento, mientras que la de un neumático real es función del ángulo de deriva, esto es, del cociente entre esa velocidad y la de avance. La consecuencia es que la rigidez de deriva efectiva del modelo crece con la velocidad en lugar de mantenerse aproximadamente constante.

% Conviene señalar que un modelo de neumático más elaborado incorpora más no linealidades conmutativas, que es precisamente la clase de dificultad que este trabajo identificó como dominante: su implementación debería cuidar la regularización con el mismo criterio con que se cuidaron la resistencia a la rodadura y el círculo de fricción.

\subsection{Distribución longitudinal de masa}

El chasis se modela como un paralelepípedo de masa uniforme, por lo que el reparto de cargas entre ejes resulta simétrico. Es una limitación que en unos casos no es relevante (cuando el vehículo se acerca a esta distribución), pero en otros puede llegar a afectar significativamente a la fidelidad de la simulación, particularmente en su comportamiento lateral.

La corrección se podría hacer creando un nuevo modelo de chasis, que modele correctamente el desplazamiento del centro de masa, y en consecuencia, de una distribución entre ejes más realista. Para esto no es necesario hacer cambios estructurales, se puede lograr creando el modelo a partir de la interfaz ya definida. Se podría tratar de experimentar utilizando \verb|BodyShape| como reemplazo de \verb|BodyBox|.

% La corrección no exige cambios estructurales. Basta con desplazar el centro de masa del chasis respecto de su centro geométrico, o componer la carrocería a partir de varios cuerpos con masas y posiciones propias. Esta última variante tiene la ventaja de representar también las inercias con más realismo, y es la que mejor se presta a caracterizar un vehículo concreto a partir de datos reales.

\subsection{Retorno de fuerza al dispositivo de entrada}

El sistema de entrada actual es unidireccional, los comandos entran al simulador, pero el conductor no recibe feedback táctil de lo que esta sucediendo dentro del entorno. Agregar un retorno de fuerza agregaría esa funcionalidad y le daría al conductor información que ahora mismo solo se puede inferir de forma visual, como perdidas de agarre o cambios de elevación bruscos.

El sistema ya tiene lo necesario para hacerlo. Las fuerzan que tienen que enviarse ya se calculan dentro del modelo y el canal de comunicación no necesita grandes cambios, solamente hay que definir un mensaje más en el protocolo y reservar un puerto extra para eso.



% Esta línea vuelve exigible una medición que este trabajo no realizó: la latencia total del lazo tal como la percibe el conductor. Está caracterizada la contribución del modelo y de su acoplamiento, pero no la del periférico ni la del renderizado. Mientras el lazo es visual, un retardo adicional pasa relativamente inadvertido; con retorno de fuerza deja de hacerlo, porque el conductor lo percibe en las manos y la tolerancia es mucho menor.

\subsection{Suavizado de la reconstrucción del comando}

Entre las mejoras de rendimiento, una mejora potencial no está en el modelo si no en la entrada que le llega. Los comandos del conductor se reciben en intervalos discretos y se mantienen constantes entre actualizaciones, por lo que cada actualización genera un escalón que el integrador debe procesar, reduciendo el paso. Reconstruir el comando de forma continua (interpolando entre muestras por ejemplo) eliminaría estas discontinuidades sin alterar la física.

% Entre las mejoras de rendimiento, la de mayor efecto potencial no está en el modelo sino en su entrada. Los comandos del conductor llegan a intervalos discretos y se mantienen constantes entre actualizaciones, de modo que cada actualización introduce un escalón que el integrador debe absorber reduciendo el paso. Reconstruir el comando de forma continua —interpolando entre muestras, o filtrando la señal con una constante de tiempo pequeña frente a la respuesta del vehículo pero grande frente al paso de comunicación— eliminaría esas discontinuidades sin alterar la física.

% Es la aplicación directa de la lección que atraviesa todo el capítulo de experimentos, y la que mejor relación ofrece entre esfuerzo de implementación y margen de cómputo recuperado.

\subsection{Unificación del entorno de ejecución}

Toda la arquitectura de co-simulación adoptada existe porque OpenModelica no permite integrar con CVODE un modelo que tenga un reloj de muestreo. Corregir este defecto en el propio compilador permitiría ejecutar la totalidad del sistema desde OpenModelica, unificando todas las herramientas de la simulación en sí, lo cual es posible debido a la naturaleza de código abierto del compilador.



% La última línea no es sobre el simulador sino sobre la herramienta. La arquitectura de co-simulación adoptada, con el muestreo y el acompasamiento fuera del modelo, existe porque OpenModelica no admite integrar con CVODE un modelo que contenga un reloj periódico de muestreo. Corregir ese defecto en el propio compilador permitiría ejecutar el lazo cerrado íntegramente dentro de OpenModelica y disponer de una única herramienta para modelar, ensayar y ejecutar, en lugar de las dos que este trabajo debió combinar.

% El defecto está documentado en el apéndice con su caso mínimo de reproducción, que es el insumo que un informe al proyecto necesita. Vale aclarar que la corrección no invalidaría la arquitectura desarrollada: el desacoplamiento entre el motor físico, la entrada y la visualización seguiría siendo deseable por sí mismo, y de hecho es lo que permite sustituir el motor físico por otro. Lo que la corrección aportaría es la posibilidad de elegir, en lugar de la obligación de rodear una limitación.
